\section{Limitation and Future Work}



\begin{table}
  \begin{center}
    \caption{Comparison of different methods on TinyImageNet splits in class-incremental settings with 100 base classes.}
    \label{tab:tiny_imagenet_class_incre}
    \resizebox{\columnwidth}{!}{%
    \begin{tabular}{y{55}{c}c|cc|cc} 
  \toprule
    & \multicolumn{2}{c}{{5 steps}} &  \multicolumn{2}{c}{10 steps}    &  \multicolumn{2}{c}{20 steps}  \\
    \textbf{Methods}    & \textbf{Avg}   & \textbf{Last}    & \textbf{Avg}   & \textbf{Last}    & \textbf{Avg}   & \textbf{Last}  \\
    \midrule
    EWC~\cite{kirkpatrick2017overcoming}       & 19.01 & 6.00  & 15.82 & 3.79  & 12.35 & 4.73 \\
    % LwF~\cite{li2017learning}                & 22.31 & 7.34  & 17.34 & 4.73  & 12.48 & 4.26 \\
    % iCaRL~\cite{rebuffi2017icarl}          & 34.27 & 23.22 & 30.94 & 20.82 & 27.83 & 20.16 \\
    EEIL~\cite{castro2018end}          & 47.17 & 35.12 & 45.03 & 34.64 & 40.41 & 29.72 \\
    UCIR~\cite{hou2019learning}               & 50.30 & 39.42 & 48.58 & 37.29 & 42.84 & 30.85 \\
    MUC~\cite{liu2020more}                    & 32.23 & 19.20  & 26.67 & 15.33 & 21.89 & 10.32 \\
    PASS~\cite{zhu2021prototype}        & 49.54 & 41.64 & 47.19 & 39.27 & 42.01 & 32.93 \\
    DyTox~\cite{douillard2022dytox}           & 55.58 & 47.23 & 52.26 & 42.79 & 46.18 & 36.21 \\
    \midrule
    CLIP~\cite{radford2021learning}  & \underline{69.62} & \underline{65.30} & \underline{69.55} & \underline{65.59} & \underline{69.49} & \underline{65.30} \\ 
    FT & 61.54 &46.66  &57.05 &41.54 & 54.62 & 44.55\\
    LwF~\cite{li2017learning} &60.97 &48.77 &57.60 &44.00 &54.79 &42.26 \\
    iCaRL~\cite{rebuffi2017icarl} & 77.02 & 70.39 & 73.48 & 65.97 & 69.65 & 64.68 \\
    LwF-VR~\cite{ding2022don} & 77.56 & 70.89 & 74.12 & 67.05 & 69.94 & 63.89 \\
    ZSCL (Ours) & \textbf{80.27} & \textbf{73.57} & \textbf{78.61} & \textbf{71.62} & \textbf{77.18} & \textbf{68.30} \\
    \textbf{Impr} & \hspace{-4pt}\gre{\textbf{+10.65}} & \gre{\textbf{+8.27}} & \gre{\textbf{+9.06}} & \gre{\textbf{+6.03}} & \gre{\textbf{+7.69}} & \gre{\textbf{+3.00}} \\
    \bottomrule
    \end{tabular}}
\end{center}
\end{table}


% generated dataset
Our proposed method has a limitation that a reference dataset is needed. A promising direction of the work is to preserve the zero-shot transfer ability without the need for an outside dataset. For example, we may generate a synthetic image dataset as the reference dataset. Methods like~\cite{smith2021always} can synthesize datasets from a network.

The deep learning community tends to build large models with a huge dataset~\cite{brown2020language,dehghani2023scaling}, including vision-language models~\cite{radford2021learning,huang2023language}. 
% These models serve as foundation models~\cite{bommasani2021opportunities} for downstream tasks and require millions of dollars for training. 
As re-training cost upsurges, continual learning is an efficient approach for updating these models with custom usage.

% correct and forget
% Our continual learning setting learns new knowledge for the vision-language model. 
In some cases, we want to correct wrong information in the pre-training dataset or update old information with the latest one. How to conduct this task with a reference dataset is left as future work.

Lately, a noticeable trend involves the creation of multi-modality models utilizing large language models, showcasing encouraging outcomes in tasks such as Visual Question Answering (VQA). Expanding our approach to encompass the next token prediction task remains an avenue for future exploration and research.

\section{Conclusion}
\label{sec:conclusion}

In this paper, we investigate continual learning with the vision-language model. We propose a better continual learning algorithm to protect the zero-shot transfer ability in the vision-language model learned in the pre-training stage. Our algorithm mitigates the catastrophic forgetting in both feature space and parameter space. In feature space, distilling the initial model on a reference dataset significantly boosts the model's performance. In parameter space, weight ensemble among different training stages alleviates the forgetting issue. We propose a more challenging Multi-domain Task Incremental Learning (MTIL) benchmark to evaluate the continual learning methods better. On both conventional and new benchmarks, our method achieves state-of-the-art performance.